\documentclass[11pt]{article}
% packages
%% font, back-end
\usepackage[english]{babel} % more, non-english characters
\usepackage{csquotes} % back-end quotation forming
\usepackage[T1]{fontenc} % back-end for non-english characters
\usepackage{graphicx} % advanced graphics support
%% formatting
\usepackage{comment} % adds the comment environment
\usepackage[letterpaper,top=2.54cm,bottom=2.54cm,left=2.54cm,right=2.54cm,marginparwidth=1.75cm]{geometry} % page dimensions
\usepackage[colorlinks=true, allcolors=blue]{hyperref} % makes hyper-links
\usepackage{ragged2e} % text alignment
\usepackage{setspace} % linespacing
\onehalfspacing

%% math, symbols
\usepackage{amsmath} % math formulas
\usepackage{array} % arrays inside math mode
\usepackage{textcomp, gensymb} % more symbols and get rid of pesky warnings
\usepackage{pdfpages}

%% figures, tables
\usepackage{booktabs} % better tables
\usepackage{longtable} % make tables that span multiple pages
\usepackage{multirow} % multi-columns and rows for tables
\usepackage{subfig} % allows for subfigures

\makeatletter
  \renewcommand\l@section{\@dottedtocline{2}{1em}{2em}}
  \renewcommand\l@subsection{\@dottedtocline{2}{2em}{3em}}
\makeatother

%% mutliple columns for the text
\usepackage{multicol}
\usepackage{blindtext}

\newcommand{\iso}[2]{$^{#2}$#1}
\renewcommand{\arraystretch}{1.5}
\setlength{\parindent}{0pt}

\begin{document}

\begin{center}
{\LARGE {PHYS 427}} \\ \vspace{5pt}
{\LARGE {Thermal and Statistical Physics}} \\ \vspace{5pt}
\end{center}

\tableofcontents
\clearpage

\section{Laws of Thermodynamics}
The zeroth law states if two systems are both in thermal equilibrium with a third system, they are in thermal equilibrium with each other. \\

The first law states the change in internal energy of a system is equal to the difference between the heat supplied to the system and the work done by the system on its surroundings.
\begin{equation}
    \Delta U = Q - W
\end{equation} \\

The second law states the change in entropy must be greater than or equal to zero. \\

The third law states the entropy of a system approaches 0 at absolute zero.


\section{Counting Microstates}
Use the choose operator
\newcommand{\Nchoosen}{\begin{equation}
    \Omega(n, N) = \frac{N!}{n! (N - n)!} \equiv \binom{N}{n}
\end{equation}}
\Nchoosen

Variance of $n$:
\begin{equation}
    \sigma_n^2 = \langle n^2 \rangle - \langle n \rangle ^2
\end{equation}

Entropy
\begin{equation}
    \sigma = \ln \Omega
\end{equation}

\begin{equation}
    \sigma_{a + b} = \sigma_a + \sigma_b 
\end{equation}

Einstein Solid
\begin{equation}
    \sigma = N \ln U
\end{equation}


\section{Entropy}
Fundamental Temperature
\begin{equation}
    \frac{1}{\tau} \equiv \frac{\partial \sigma}{\partial U} \equiv \frac{1}{kT}
\end{equation}

Conventional Entropy
\begin{equation}
    S = k \ln \Omega
\end{equation}

\begin{equation}
    \frac{1}{T} \equiv \frac{\partial S}{\partial U}
\end{equation}

Pressures are equal at mechanical equilibrium
\begin{equation}
    p = T \left( \frac{\partial S}{\partial V} \right)_{U,N}
\end{equation}


\section{Canonical Ensemble}
For all possible energies, the partition function is defined as:
\begin{equation}
    Z = \Sigma_i e^{-\epsilon_i/kT} = \Sigma_i e^{-\beta \epsilon_i}
\end{equation}

Each state has a probability given by the ratio of the Boltzmann factor and the partition function:
\begin{equation}
    p = \frac{\Sigma_i e^{-\beta \epsilon_i}}{Z}
\end{equation}

Energy given by Z:
\begin{equation}
    U = \langle \varepsilon \rangle = - \frac{\partial}{\partial \beta} \ln Z
\end{equation}

For a paramagnet
\begin{equation}
    U = -\mu B \tanh{\beta \mu B}
\end{equation}

Equipartition theorem states that the partition function of a system of $N$ particles is the product of the individual partition functions. If the particles are identical and distinguishable, the partition function is given as:
\begin{equation}
    Z_{total} = Z_1^N
\end{equation}

If the particles are indistinguishable,
\begin{equation}
    Z = \frac{Z_1^N}{N!}
\end{equation}


\subsection{Quantum Density}
The quantum density is defined such that
\begin{equation}
    Z = n_Q V
\end{equation}

Therefore, the quantum density is given as
\begin{equation}
    n_Q \equiv \left( \frac{mkT}{2\pi \hbar^2} \right)^{3/2}
\end{equation}

If the density $n$ is much less than $n_Q$, then this is called the classical regime. In the classical regime, quantum effects are unimportant.


\subsection{Equipartition Theorem}
For many high-temperature cases, $U = \alpha NkT$ (Einstein model, ideal gas, diatomic gas) where $alpha$ is the number of quadratic degrees of freedom.

For a Maxwell distribution,
\begin{align}
    v_{peak} = \sqrt{\frac{2kT}{m}} \\
    \langle v \rangle = \sqrt{\frac{8kt}{\pi m}} \\
    \langle v^2 \rangle = \frac{3kt}{m}
\end{align}

\subsection{Sackur-Tetrode Equation}
Equation that gives the entropy of an ideal gas
\begin{equation}
    S = Nk \left[ \ln{\frac{n_QV}{N}} + \frac{5}{2} \right]
\end{equation}
\begin{equation}
    F = - NkT \left[ \ln{\frac{n_QV}{N}} + 1 \right]
\end{equation}


\section{Thermodynamic Potentials}

\begin{itemize}
    \item Helmholtz is the work provided to create a system at constant T.
    \item Enthalpy is the work required to create a system and displace the gas in the volume of the system.
    \item Gibbs is enthalpy minus T$\Delta$S because the heat of the gas was also removed.
\end{itemize}

Cannot have potentials with natural variables of V,P or T,S because these are conjugate pairs.
\begin{center}
\begin{tabular}{c|c|c|c}
    {\bf Energy} & {\bf Nat. Vars.} & {\bf Expression} & {\bf Differential} \\ \hline
    Internal & S,V & $U = U(S, V)$ & $dU = TdS - pdV$ \\
    Helmholtz & T,V & $F \equiv U - TS$ & $dF = -SdT - pdV$ \\
    Enthalpy & S,P & $H \equiv U + pV$ & $dH = TdS + Vdp$ \\
    Gibbs & T,P & $G \equiv U + pV - TS$ & $dG = Vdp - SdT$
\end{tabular}
\end{center}

\begin{equation}
    F = - kT \ln{Z}
\end{equation}

\subsection{Chemical Potential}
If N is allowed to change,
\begin{equation}
    \mu = 
    -T \left( \frac{\partial S}{\partial N} \right)_{U,V} = 
    \left( \frac{\partial U}{\partial N} \right)_{S,V} = 
    \left( \frac{\partial F}{\partial N} \right)_{T,V}
\end{equation}

For some external force, $\mu$ takes on the energy associated with that force. For gravity,
\begin{equation}
    \mu = kT \ln{\frac{n(h)}{n_Q(T)}} + mgh
\end{equation}

For an ideal gas, and generally,
\begin{equation}
    \mu = kT \ln{\frac{n}{n_Q}}
\end{equation}

When multiple species are in equilibrium:
\begin{equation}
    \mu_1 + \mu_2 + ... + \mu_n = 0
\end{equation}


\section{Grand Canonical Ensemble}
Gibbs Factor
\begin{equation}
    e^{-(\varepsilon_j - \mu N_i) / kT}
\end{equation}

Grand partition function:
\begin{equation}
    \xi = \sum_{N_i} \sum_{\varepsilon_j} e^{-(\varepsilon_j - \mu N_i) / kT}
\end{equation}
\begin{equation}
    \Xi = \sum_{N_i} \lambda^N Z(N, T)
\end{equation}
With $\lambda \equiv e^{\mu / kT}$ as the absolute activity.



\section{Math Formulae}
Choose Operator
\Nchoosen

Sterling Approximation
\begin{equation}
    \ln N ! \simeq N \ln N - N
\end{equation}

\end{document}